\begin{theorem}[S-30: An integral lattice obstruction closes the simple rank-one balanced RM gate]
\label{thm:semireg-rm-integral-obstruction}
Let $(X,\theta)$ be a complex abelian fourfold with a
Rosati-invariant action of the real quadratic field
$F=\mathbf Q(\sqrt{\Delta_0})$, where $\Delta_0>1$ is
squarefree, and with the two complex eigenspaces both
of dimension two. Use the notation
\[
 D=\int_X\frac{\theta^4}{24},\qquad
 C=\frac{\theta_1^2+\theta_2^2}{2},\qquad
 v_q=1-qC+\frac{q^2}{4}\theta_1^2\theta_2^2,
 \qquad q\in\mathbf Q_{>0}.
\]
Here $\theta$ is an integral polarization and
$\theta=\theta_1+\theta_2$ is its orthogonal decomposition
along the two eigenspaces.
If $a\in\mathbf Q$ is nonzero and $aC$ is an integral
cohomology class, then
\[
 t=a^2D\in\mathbf Z_{>0},\qquad \Delta_0\mid t.
 \tag{S30.1}
\]
Consequently no rank-one perfect object with normalized
character $v_q$ can be both simple, gluable, and fully
semiregular. More quantitatively, every simple gluable
rank-one perfect object with that normalized character
satisfies
\[
 \dim\ker\sigma_E^2\ge2Dq^2-2\ge2\Delta_0-2>0.
 \tag{S30.2}
\]
In particular, in the principal $q=1$ case, the character
$v_1$ is not the normalized character of any rank-one
perfect object at all. This last assertion does not
require simplicity, gluability, or semiregularity.

For a positive-rank perfect object with
$e^B\operatorname{ch}(E)=Nv_q$, one always has
\[
 N^4Dq^2\in\mathbf Z_{>0},\qquad
 \Delta_0\mid N^4Dq^2.
 \tag{S30.3}
\]
If, in addition, $E$ has a simple gluable fully semiregular
lift for an effective finite coherent target action $G$
as in S-22, write $k=\chi_G\in\{2,4\}$. Then
\[
 \begin{cases}
 \Delta_0\mid N^2|G|,&k=4,\\
 2\Delta_0\mid N^2|G|,&k=2.
 \end{cases}
 \tag{S30.4}
\]
For $G=1$, in particular, $N$ must be a multiple of
$\Delta_0$, and $N$ must also be even in the $k=2$ case.
Moreover, in every rank, a nonequivariant simple gluable
fully semiregular realization requires a genuinely
fractional normalizing class: $B$ cannot be the Chern
class of a line bundle. Equivalently, its determinant
cannot admit an $N$th root in $\operatorname{Pic}(X)$.
The integer $\Delta_0$ here is the squarefree radicand,
not necessarily the discriminant of the quadratic field.
\end{theorem}
\begin{proof}
Put $L=H_1(X,\mathbf Z)$ and let $\nu$ be the primitive
positive generator of $\bigwedge^8L^*$, so
$\theta^4/24=D\nu$. The two real eigenspaces give
$L_{\mathbf R}=U\oplus W$, each of real dimension four.
Write $\nu_U=\theta_1^2/2$ and $\nu_W=\theta_2^2/2$.
Then $C=\nu_U+\nu_W$ and $\nu_U\wedge\nu_W=D\nu$.

Suppose $c=aC\in\bigwedge^4L^*$. Define the integral
endomorphism $R_c$ of $\bigwedge^2L$ by
\[
 \iota_{R_c(x)}\nu=(\iota_xc)\wedge c.
 \tag{S30.5}
\]
Contraction with the primitive volume form is a
unimodular identification between
$\bigwedge^2L$ and $\bigwedge^6L^*$, so $R_c$ is indeed
integral. A direct calculation on the three summands
\[
 \bigwedge^2L_{\mathbf R}
  =\bigwedge^2U\oplus(U\otimes W)\oplus\bigwedge^2W
\]
shows that $R_c=tP$, where $t=a^2D$ and $P$ is the
projector which is the identity on the first and third
summands and zero on the middle summand.
Moreover $c^2=2t\nu$. The square of every integral
four-form on an eight-dimensional lattice has even top
coefficient: basis monomials square to zero and their
cross terms occur in pairs. Hence $t$ is an integer.

Choose the rational endomorphism $T$ of $L_{\mathbf Q}$
given by $\sqrt{\Delta_0}\in F$. It satisfies
$T^2=\Delta_0\operatorname{id}$ and has eigenvalues
$\sqrt{\Delta_0},-\sqrt{\Delta_0}$ on $U,W$. Therefore
\[
 2P-\operatorname{id}
      =\frac{\bigwedge^2T}{\Delta_0}.
 \tag{S30.6}
\]
Let $p$ divide $\Delta_0$. If $p\nmid t$, the integrality
of $tP$ makes $P$ integral over $\mathbf Z_p$.
Thus the left side of (S30.6) is an integral involution,
and hence an automorphism of $\bigwedge^2L_{\mathbf Z_p}$.
Apply Smith normal form over $\mathbf Z_p$ to the
invertible rational matrix $T$. If its elementary-divisor
exponents are integers $m_1,\ldots,m_8$, those of its
exterior square, before reordering, are
$m_i+m_j$ for $i<j$. Equation (S30.6) and unimodularity
would force
\[
 m_i+m_j=v_p(\Delta_0)=1\qquad(i<j).
 \tag{S30.7}
\]
Using any three indices forces every $m_i=1/2$, a
contradiction. This argument works also at $p=2$.
Every prime factor of the squarefree integer $\Delta_0$
therefore divides $t$, proving (S30.1).

For rank one, the degree-two character equation gives
$B=-c_1(E)$. The determinant of a perfect complex is a
line bundle, so $E'=E\otimes\det(E)^{-1}$ is an actual
perfect complex with character $v_q$. Its second Chern
class is $qC$, which is integral. Riemann--Roch on the
abelian fourfold gives $\chi(E')=Dq^2$ and
$\chi(E,E)=4Dq^2$. By (S30.1), $Dq^2$ is a positive
multiple of $\Delta_0$.

If $E$ is simple and gluable, the rank-eight degree-one
contraction forces $e_1\ge8$, and CY4 Serre duality gives
\[
 m:=e_2=4Dq^2-2+2e_1\ge4Dq^2+14.
\]
Put $b=\dim\ker\sigma_E^2$. The full transpose identity
makes the ranks of $\sigma_E^2$ and
$\operatorname{ob}_E^2$ both $m-b$. Their composition has
rank eighteen. Thus
$18\ge2(m-b)-m=m-2b$, which proves (S30.2).
For $D=q=1$, (S30.1) already contradicts $qC$ being
integral; this proves the assertion without simplicity.

For general positive rank $N$, put $\ell=c_1(E)$.
The degree-two equation gives $B=-\ell/N$, and the
next character equation is
\[
 N^2qC=Nc_2(E)-(N-1)\frac{\ell^2}{2}.
 \tag{S30.8}
\]
The class $\ell^2/2$ is integral in the exterior algebra
of a torus. Thus the right side is integral and (S30.1)
with $a=N^2q$ proves (S30.3). In the equivariant case
S-22 gives $k=4Dq^2N^2/|G|\in\{2,4\}$, so
\[
 N^4Dq^2=\frac{kN^2|G|}{4}.
\]
Substitution into (S30.3) gives (S30.4).
For $G=1$, squarefreeness of $\Delta_0$ implies
$\Delta_0\mid N$; when $k=2$, the additional divisibility
by two forces $N$ even.

Finally suppose $G=1$ and $B=c_1(Q)$ for a line bundle
$Q$. Then $E\otimes Q$ is an actual perfect object with
character $Nv_q$, so $s:=\chi(E\otimes Q)=NDq^2$ is a
positive integer. Full semiregularity, simplicity and
gluability give
$4Ns=\chi(E,E)=2-2e_1+18\le4$.
Thus $N=s=1$, which is the excluded rank-one case.
Since $NB=-c_1(E)$, a line-bundle determinant root would
make $B$ integral, and conversely an integral $B$ is a
line-bundle class by the integral $(1,1)$ theorem.
The residual degree-zero line bundle can be absorbed
into an $N$th root because multiplication by $N$ on
$\operatorname{Pic}^0(X)$ is surjective.
\end{proof}

\begin{proposition}[S-31: Rank-two and rank-three bundles cannot realize the balanced character]
\label{prop:semireg-rm-low-rank-bundles}
In the real-multiplication setting of S-30, no vector
bundle of rank $N\in\{1,2,3\}$ has normalized character
$e^B\operatorname{ch}(E)=Nv_q$ with $q>0$.
Consequently a nonequivariant simple gluable fully
semiregular vector-bundle realization would have rank
at least four and divisible by $\Delta_0$.
For example, the smallest remaining ranks allowed by
these two tests are four for $F=\mathbf Q(\sqrt2)$,
and six for $F=\mathbf Q(\sqrt3)$.
These are necessary conditions, not constructions.
\end{proposition}
\begin{proof}
Under the splitting principle, replace the $N$ Chern
roots $x_i$ by $x_i+B$. The resulting normalized Chern
polynomial still has degree at most $N$, regardless of
whether $B$ is the Chern class of a line bundle.
The normalized character has first and third components
zero, second component $-NqC$, and fourth component
$NDq^2\nu$, with $\nu$ the primitive top class.
Newton's identities therefore give
\[
 \bar c_1=\bar c_3=0,\qquad
 \bar c_2=NqC,\qquad
 \bar c_4=\frac{\bar c_2^2-12\overline{\operatorname{ch}}_4}{2}
          =N(N-6)Dq^2\nu.
 \tag{S31.1}
\]
For $N=1$, the nonzero second Chern class contradicts
the degree-one Chern polynomial. For $N=2,3$, the
nonzero fourth class contradicts degree at most $N$.
The final rank restrictions follow from S-30.
\end{proof}
